\documentclass[11pt,a4paper]{article}
\usepackage[margin=25mm]{geometry}
\usepackage{amsmath,amssymb,amsthm,hyperref}
\newtheorem{theorem}{Theorem}
\newtheorem{lemma}[theorem]{Lemma}
\newcommand{\gen}[1]{\langle#1\rangle}
\newcommand{\F}{\mathbf F}
\title{The four displayed symplectic generators}
\author{Human Compression project: L031 corrected recovery}
\date{October 2026}
\begin{document}\maketitle
We verify the four matrices printed after Theorem 6.1 of
Brooksbank--Ferrara--Leemans, in the supplied twelve-page edition
of arXiv:1709.02219v2. They do give a rank-four string C-marking
for every stated parameter. Their first three generators preserve
a \emph{plus}-type quadratic form; thus the preceding minus-type
identification is not the identification of this displayed family.
The proof below retains the complete ambient support and three-space
lemmas of the companion \emph{Endpoint variation in orthogonal
symmetry markings}. It does not retain a finite experiment or an
orthogonal maximal-subgroup assertion.

\section{Parameters, form and two dihedral groups}
Let $F=\F_q$, $q\geq4$ even, and choose $\alpha\in F^\times$ such that
\[
 C_\alpha=\begin{pmatrix}0&1\\1&\alpha^2\end{pmatrix}
 \quad\hbox{has order }q+1.
\]
Such parameters exist by the norm-one torus construction in the
companion. In particular $\alpha\ne1$, since $C_1$ has order three.
Vectors are rows in the ordered basis $(u,e,b,w)$. Define
\[
 r_0=\begin{pmatrix}1&0&0&0\\\alpha&1&0&0\\0&0&1&0\\0&0&0&1\end{pmatrix},
 \quad r_1=\begin{pmatrix}1&\alpha&\alpha&0\\0&1&0&0\\0&0&1&0\\0&1+\alpha&\alpha&1\end{pmatrix},
\]
\[
 r_2=\begin{pmatrix}1&0&0&0\\0&1&0&\alpha\\0&0&1&1\\0&0&0&1\end{pmatrix},
 \quad r_3=\begin{pmatrix}1&0&0&0\\0&1&0&0\\0&0&1&0\\0&0&\alpha&1\end{pmatrix}.
                                                               \tag{1}
\]
Use the quadratic form
\[
 Q(U,E,B,W)=U^2+B^2+W^2+\alpha UE+\alpha EW+BW.       \tag{2}
\]
Its polar matrix is
\[
 \begin{pmatrix}0&\alpha&0&0\\\alpha&0&0&\alpha\\
 0&0&0&1\\0&\alpha&1&0\end{pmatrix},
\]
which is nondegenerate. We write its pairing as $(\ ,\ )$ and
$s_v(x)=x+(x,v)v/Q(v)$ for its symmetries. Direct substitution gives
$r_0=s_u$, $r_2=s_w$, and $r_3(x)=x+\alpha(x,b)b$.
Put
\[
 l=e+b,\quad h=(1+\alpha)\alpha^{-1}e+b,\quad y=e+\alpha b,
 \qquad L_u=\gen{u,l},\quad L_w=\gen{w,h}.
\]
The values $Q(u),Q(l),Q(w),Q(h),Q(b)$ are one, while $Q(y)=\alpha^2$.
Moreover $(u,l)=(w,h)=\alpha$, and
\[
                 r_1=s_ls_b=s_hs_y.                         \tag{3}
\]
Both identities follow by applying the two sides to $u,e,b,w$;
their two symmetry factors commute. In particular the first three
matrices preserve $Q$. All four preserve its polar form and are
nonidentity involutions. The nonadjacent pairs in (1) commute.

The characteristic polynomial of $s_us_l$ on $L_u$, and of $s_ws_h$
on $L_w$, is $X^2+\alpha^2X+1$. These elements have order $q+1$,
by the condition on $C_\alpha$. Hence $L_u,L_w$ are asingular and
their displayed symmetry pairs generate their whole dihedral
isometry groups. Since $b\perp L_u$ and $y\perp L_w$, (3) and the
oddness of $q+1$ give
\[
 \gen{r_0,r_1}=O(L_u)\times\gen{s_b},\qquad
 \gen{r_1,r_2}=O(L_w)\times\gen{s_y}.                         \tag{4}
\]
These are actual ambient internal direct products, of order $4(q+1)$:
for example $(r_0r_1)^{q+1}=s_b$. Both factors fix the perpendicular
complement of their stated support.

The type of (2) is explicitly plus. Indeed
$v=w+(1+\alpha^{-1})u$ lies in $L_u^\perp$ and
\[
 Q(v)=\alpha^{-2},\quad Q(b)=1,\quad(v,b)=1.
\]
The two perpendicular planes $L_u$ and $\gen{b,v}$ have the same
Arf trace $\operatorname{Tr}(\alpha^{-2})=1$. Their sum therefore
has Arf trace zero. In particular the distinction from the source's
minus-type discussion is a uniform calculation, not an inference
from a sampled order.

\section{Full generation}
Let $H=\gen{r_0,r_1,r_2}$. Equation (4) supplies $\Sigma_{L_u}$
and $s_b$. The space $T=\gen{u,l,w}$ is orthogonal of dimension
three. Its radical is generated by $(\alpha+1)u+\alpha w$, of
quadratic value one. Also $w$ is nonradical in $T$, since
$(w,l)=\alpha+1\ne0$. The companion's ambient three-space lemma
gives $\Sigma_T\leq H$.

The asingular plane $L_u^{s_w}$ contains
$l+(\alpha+1)w$, whose pairing with $b$ is $\alpha+1\ne0$.
Consequently its conjugate by $s_b$ is an asingular plane not
contained in $T$: the changed vector acquires the nonzero
$b$-component $(\alpha+1)b$ modulo $T$.
The companion's hyperplane saturation lemma now gives
\[
                         H=O(V,Q)=O^+(4,q).                \tag{5}
\]

To see that $G=\gen{H,r_3}$ is the entire symplectic group without
using maximality, choose $t\ne0$ with
$t^2=(\alpha^2+\alpha)^{-1}$ and set
\[
       a=(0,(1+t^2)/(\alpha t),0,t).
\]
Then $Q(a)=1$, $(a,b)=t$, and
\[
 Q(a r_3)=Q(a)+(\alpha^2+\alpha)(a,b)^2=0.
\]
Thus conjugating $s_a\in H$ by $r_3$ gives a transvection with
a nonzero singular axis. The full plus orthogonal group is transitive
on its nonzero singular vectors: in the hyperbolic model
$Q(X)=\det X$ on $2$ by $2$ matrices, these are the nonzero
rank-one matrices, and left and right $\mathrm{SL}_2(q)$ actions
are transitive on them. Scaling an axis scales its transvection
parameter by a square; all nonzero parameters therefore occur
on every singular axis.

These transvections generate every transvection. For a nonsingular
$z$, choose a singular $x$ with $(z,x)\ne0$; singular vectors span
the hyperbolic space, so this is possible. The vector
$z'=z+Q(z)x/(z,x)$ is singular. A transvection on $x$ sends $z'$
to $z$. Conjugation consequently supplies every parameter on $z$.
The elementary transvection generation in the companion proves
\[
                            G=\operatorname{Sp}(4,q).       \tag{6}
\]

\section{All intersections}
First consider the three marks of $H$. Put
$U=\gen{u,e,b}$ and $W=\gen{w,e,b}$. Elements of the two
groups in (4) have support in $U$ and $W$ respectively; their
intersection is supported in $\gen{e,b}$.
An element of the first group supported there preserves $L_u$
and has its image on $L_u$ contained in $\gen l$. The
one-space support assertion in the companion makes its action
on $L_u$ either $1$ or $s_l$. Its complementary factor is
$1$ or $s_b$. Thus it belongs to $\gen{s_l,s_b}$.
Likewise the second group contributes only $\gen{s_h,s_y}$.
The root lines $\gen l,\gen b,\gen h,\gen y$ are pairwise
distinct when $\alpha\notin\{0,1\}$. The first two Klein four
groups have, by (3), their sole common nonidentity element $r_1$.
Their intersection is therefore exactly $\gen{r_1}$.
This proves that $(r_0,r_1,r_2)$ is C.

For the other end group $K=\gen{r_1,r_2,r_3}$, the second group
in (4) preserves $L_w$ and its perpendicular complement. If one
of its elements is also in $\gen{r_2,r_3}$, its support lies in
$\gen{w,b}$. On $L_w$ it is supported on
$L_w\cap\gen{w,b}=\gen w$, so is $1$ or $s_w$ there.
Its possible other factor is $s_y$, whose root is outside
$\gen{w,b}$. Hence the intersection is exactly $\gen{r_2}$.
Both adjacent two-generator groups consist of distinct involutions,
so this proves that $(r_1,r_2,r_3)$ is C.

It remains to intersect $H$ and $K$. Every element of $K$ is
supported in $W$, and preserves on $W$ the quadratic form
\[
 Q'(E,B,W)=\frac{\alpha}{\alpha+1}E^2+
                 \alpha^{-1}B^2+W^2+\alpha EW+BW.            \tag{7}
\]
This is checked directly on its three generators. The difference
$Q'|_W-Q|_W$ is a nonzero squared linear functional whose kernel
is $L_w$; explicitly it vanishes on $w,h$ and is nonzero on $y$.
Thus an element of $H\cap K$ preserves $L_w$ and $Q$, and has
support in $W$. On the perpendicular decomposition
$V=L_w\perp L_w^\perp$, its first action belongs to $O(L_w)$.
Its second is supported in $W\cap L_w^\perp=\gen y$, and is
therefore $1$ or $s_y$. By (4) this yields
\[
                  H\cap K=\gen{r_1,r_2}.                  \tag{8}
\]
The standard string end-parabolic criterion proves the full
intersection property for (1).

\begin{theorem}
For every even $q\geq4$ and every source parameter $\alpha$
with $|C_\alpha|=q+1$, the four literal matrices (1) form a
rank-four string C-marking of $\operatorname{Sp}(4,q)$.
Their first-three subgroup is $O^+(4,q)$, of order
$2q^2(q^2-1)^2$, rather than the minus group of order
$2q^2(q^4-1)$.
\end{theorem}
The marking is irreducible: a commuting adjacent pair in a string
C-marking would split the generated group into two nontrivial
direct factors, whereas $\operatorname{Sp}(4,q)$ is simple for
$q\geq4$. Alternatively noncommutation is immediate in (1).
The source's principal existence conclusion is preserved. This
calculation does not assert that every preliminary geometric choice
in its preceding paragraph has the same type, and does not claim
an historical new construction.
\end{document}
